\documentclass[11pt,a4paper]{article}

\usepackage[utf8]{inputenc}
\usepackage[T1]{fontenc}
\usepackage[french]{babel}
\usepackage[a4paper,margin=2cm]{geometry}
\usepackage{amsmath,amssymb,siunitx}
\usepackage{lmodern}
\usepackage{enumitem}
\usepackage{fancyhdr}
\usepackage{lastpage}
\usepackage{hyperref}
\usepackage{xcolor}
\usepackage{array}
\usepackage{tabularx}
\usepackage{graphicx}
\usepackage{multicol}
\usepackage{bm}

\sisetup{locale=FR}

\setlength{\headheight}{14pt}
\pagestyle{fancy}
\fancyhf{}
\fancyhead[L]{Terminale générale -- Spécialité Physique-Chimie}
\fancyhead[C]{DS inspiré de GEIPI Polytech -- version approfondie}
\fancyhead[R]{Page \thepage/\pageref{LastPage}}
\fancyfoot[C]{\textit{On valorise le sens physique, la modélisation et l’esprit critique}}

\begin{document}

\begin{center}
    {\Large \textbf{DEVOIR SURVEILLÉ -- PHYSIQUE}}\\[0.3em]
    {\large \textbf{Sujet inspiré des concours GEIPI Polytech}}\\[0.3em]
    {\large \textbf{Version difficile, avec ouvertures hors programme guidées}}\\[0.5em]
    Durée conseillée : 3 h
\end{center}

\vspace{0.4em}

\noindent\textbf{Consignes}
\begin{itemize}[leftmargin=1.8em]
    \item Les trois exercices sont indépendants.
    \item Toute réponse doit être justifiée.
    \item Certaines notions non strictement au programme sont introduites dans l’énoncé : il faut les exploiter intelligemment.
    \item Une attention particulière sera portée à la cohérence physique des résultats.
\end{itemize}

\vspace{0.4em}

\begin{center}
\begin{tabular}{|>{\centering\arraybackslash}m{3.2cm}|>{\centering\arraybackslash}m{3.2cm}|>{\centering\arraybackslash}m{3.2cm}|}
\hline
Exercice 1 & Exercice 2 & Exercice 3 \\
\hline
12 points & 14 points & 18 points \\
\hline
\end{tabular}
\end{center}

\vspace{0.5em}
\hrule
\vspace{0.6em}

\section*{Exercice 1 -- Optique ondulatoire : mesurer, interpréter, critiquer \hfill /12}

On cherche à caractériser une membrane percée de trous microscopiques régulièrement espacés.  
On éclaire la membrane avec un laser vert de longueur d’onde
\[
\lambda=\SI{532}{nm}.
\]
L’écran est placé à la distance
\[
D=\SI{2,50}{m}.
\]

\medskip

\noindent\textbf{Données :}
\begin{itemize}
    \item célérité de la lumière dans le vide :
    \[
    c=\SI{3,00e8}{m.s^{-1}};
    \]
    \item pour un trou circulaire de diamètre $a$, on admet :
    \[
    \theta \simeq \frac{1,22\,\lambda}{a};
    \]
    \item pour deux ouvertures distantes de $b$, on admet :
    \[
    i=\frac{\lambda D}{b}.
    \]
\end{itemize}

\subsection*{Partie A -- Questions de base}
\begin{enumerate}[label=\textbf{1.\arabic*}, leftmargin=1.6em]
    \item Calculer la fréquence du laser.
    \item Quelle propriété fondamentale de la lumière est mise en évidence par la diffraction ? par les interférences ?
    \item Expliquer brièvement pourquoi un modèle purement corpusculaire de la lumière est insuffisant pour rendre compte de ces observations.
\end{enumerate}

\subsection*{Partie B -- Mesure d’un diamètre de trou}
On observe sur l’écran une tache centrale de diffraction de largeur
\[
L=\SI{8,0}{mm}.
\]

\begin{enumerate}[label=\textbf{1.\arabic*}, resume, leftmargin=1.6em]
    \item En supposant les petits angles, établir la relation
    \[
    \theta \simeq \frac{L}{2D}.
    \]
    \item En déduire l’expression de $a$ en fonction de $\lambda$, $D$ et $L$.
    \item Calculer $a$.
    \item Le résultat obtenu est-il compatible avec l’idée d’un \og trou microscopique \fg{} ? Commenter.
\end{enumerate}

\subsection*{Partie C -- Mesure d’un espacement entre trous}
On observe ensuite une figure d’interférences d’interfrange
\[
i=\SI{1,25}{cm}.
\]

\begin{enumerate}[label=\textbf{1.\arabic*}, resume, leftmargin=1.6em]
    \item Calculer la distance $b$ entre deux trous voisins.
    \item Comparer $a$ et $b$.
    \item Expliquer pourquoi la figure réellement observée peut être interprétée comme des interférences \emph{modulées} par une enveloppe de diffraction.
\end{enumerate}

\subsection*{Partie D -- Ouverture hors programme guidée}
Dans certains montages, le contraste des interférences diminue lorsque la différence de marche entre les deux ondes devient trop grande.  
On dit alors que la lumière n’est plus suffisamment \og cohérente \fg.

\begin{enumerate}[label=\textbf{1.\arabic*}, resume, leftmargin=1.6em]
    \item Sans définition mathématique compliquée, expliquer ce que peut signifier physiquement cette idée de cohérence.
    \item Pourquoi un laser est-il généralement mieux adapté qu’une lampe classique pour observer des interférences nettes ?
    \item Proposer deux causes expérimentales pouvant dégrader la netteté des figures observées.
\end{enumerate}

\vspace{0.5em}
\hrule
\vspace{0.6em}

\section*{Exercice 2 -- Refroidissement d’une boisson : du modèle simple à sa critique \hfill /14}

On étudie le refroidissement d’une boisson chaude dans une tasse.  
On note $\theta(t)$ sa température à l’instant $t$, et $\theta_{\mathrm{amb}}$ la température ambiante.

On admet dans un premier modèle que
\[
\frac{d\theta}{dt}=-k\big(\theta(t)-\theta_{\mathrm{amb}}\big)
\]
où $k$ est une constante positive.

\medskip

\noindent\textbf{Données :}
\[
\theta_{\mathrm{amb}}=\SI{20,0}{^\circ C},
\qquad
\theta(0)=\SI{85,0}{^\circ C},
\qquad
k=2,1\times 10^{-4}\ \si{s^{-1}}.
\]

\subsection*{Partie A -- Lecture physique du modèle}
\begin{enumerate}[label=\textbf{2.\arabic*}, leftmargin=1.6em]
    \item Vérifier par analyse d’unités que l’unité de $k$ est cohérente.
    \item Expliquer le signe de $\dfrac{d\theta}{dt}$ au début de l’expérience.
    \item Que se passe-t-il si $\theta(t)=\theta_{\mathrm{amb}}$ ?
    \item Pourquoi ce modèle est-il raisonnable au moins qualitativement ?
\end{enumerate}

\subsection*{Partie B -- Résolution}
On admet que la solution de cette équation différentielle est
\[
\theta(t)=\theta_{\mathrm{amb}}+\big(\theta(0)-\theta_{\mathrm{amb}}\big)e^{-kt}.
\]

\begin{enumerate}[label=\textbf{2.\arabic*}, resume, leftmargin=1.6em]
    \item Calculer la température au bout de \SI{10}{min}.
    \item Déterminer le temps nécessaire pour atteindre \SI{55}{^\circ C}.
    \item Déterminer la limite de $\theta(t)$ quand $t\to +\infty$ et l’interpréter.
\end{enumerate}

\subsection*{Partie C -- Une notion hors programme léger : constante de temps}
On appelle \emph{constante de temps} la grandeur
\[
\tau=\frac{1}{k}.
\]

\begin{enumerate}[label=\textbf{2.\arabic*}, resume, leftmargin=1.6em]
    \item Calculer $\tau$.
    \item Montrer qu’au temps $t=\tau$,
    \[
    \theta(\tau)=\theta_{\mathrm{amb}}+\frac{\theta(0)-\theta_{\mathrm{amb}}}{e}.
    \]
    \item Interpréter physiquement ce résultat.
    \item Une tasse mieux isolée conduit-elle à une valeur plus grande ou plus petite de $\tau$ ? Justifier.
\end{enumerate}

\subsection*{Partie D -- Aller plus loin : critique du modèle}
En réalité, l’évaporation à la surface du liquide peut jouer un rôle important, surtout lorsque la boisson est très chaude.

\begin{enumerate}[label=\textbf{2.\arabic*}, resume, leftmargin=1.6em]
    \item Expliquer qualitativement pourquoi l’évaporation accélère le refroidissement.
    \item Le coefficient $k$ a-t-il de bonnes chances d’être vraiment constant au cours de toute l’expérience ? Discuter.
    \item Un élève affirme : \og si on met deux fois plus de boisson dans la tasse, le temps de refroidissement double exactement \fg.  
    Discuter cette affirmation avec un raisonnement physique.
\end{enumerate}

\vspace{0.5em}
\hrule
\vspace{0.6em}

\section*{Exercice 3 -- Mouvement d’un skieur : énergie, trajectoire, frottements \hfill /18}

Un skieur descend un tremplin puis s’élance dans les airs.  
On modélise le skieur par un point matériel de masse
\[
m=\SI{60}{kg}.
\]
On prendra
\[
g=\SI{9,81}{m.s^{-2}}.
\]

\subsection*{Partie A -- Descente sans frottements}
Le skieur part du repos depuis un point $D$ situé à la hauteur
\[
h=\SI{35}{m}
\]
au-dessus du point d’envol $E$.

\begin{enumerate}[label=\textbf{3.\arabic*}, leftmargin=1.6em]
    \item En négligeant les frottements, établir l’expression de la vitesse $V_E$ au point $E$.
    \item Calculer cette vitesse.
    \item Calculer l’énergie cinétique au point $E$.
    \item Commenter l’ordre de grandeur obtenu.
\end{enumerate}

\subsection*{Partie B -- Premier niveau de réalisme}
On suppose maintenant qu’entre $D$ et $E$, des frottements dissipent une énergie mécanique totale
\[
E_f=\SI{3,5e3}{J}.
\]

\begin{enumerate}[label=\textbf{3.\arabic*}, resume, leftmargin=1.6em]
    \item Écrire le bilan énergétique adapté.
    \item Calculer la nouvelle vitesse d’envol.
    \item Comparer ce résultat à celui de la partie A.
    \item Un écart de quelques \si{m.s^{-1}} sur la vitesse d’envol peut-il avoir des conséquences importantes sur le saut ? Expliquer.
\end{enumerate}

\subsection*{Partie C -- Vol dans l’air}
Le skieur quitte le tremplin avec la vitesse
\[
V_0=\SI{22,0}{m.s^{-1}}
\]
sous l’angle
\[
\alpha=11^\circ
\]
par rapport à l’horizontale.

On néglige l’action de l’air dans cette partie.  
On choisit un repère $(Ox,Oz)$ tel que le point d’envol soit l’origine.

\begin{enumerate}[label=\textbf{3.\arabic*}, resume, leftmargin=1.6em]
    \item Établir les équations horaires $x(t)$ et $z(t)$.
    \item En déduire l’équation de la trajectoire.
    \item Déterminer la hauteur maximale atteinte au-dessus du point d’envol.
    \item Déterminer la durée du vol si le skieur retombe à la même altitude que le point d’envol.
    \item En déduire la portée horizontale.
\end{enumerate}

\subsection*{Partie D -- Ouverture hors programme guidée : frottement fluide simplifié}
On souhaite discuter l’effet de l’air.  
Dans un modèle simplifié, on suppose que pendant le vol, l’air exerce une force opposée au mouvement, de norme proportionnelle à la vitesse :
\[
f=kv
\]
avec $k>0$.

On ne demande \textbf{aucune résolution complète} du mouvement.

\begin{enumerate}[label=\textbf{3.\arabic*}, resume, leftmargin=1.6em]
    \item Pourquoi cette force tend-elle à diminuer la portée du saut ?
    \item Cette force agit-elle seulement sur la valeur de la vitesse ou aussi sur la direction du vecteur vitesse ? Expliquer.
    \item À votre avis, l’effet de l’air est-il plus marqué au début du vol ou à la fin ? Discuter.
    \item Un skieur adopte en vol une position plus aérodynamique. Expliquer qualitativement pourquoi cela peut améliorer la performance.
\end{enumerate}

\subsection*{Partie E -- Question ouverte de synthèse}
Pour améliorer un saut, on peut agir sur :
\begin{itemize}
    \item la hauteur de départ ;
    \item les frottements sur le tremplin ;
    \item l’angle d’envol ;
    \item l’aérodynamique en vol.
\end{itemize}

\begin{enumerate}[label=\textbf{3.\arabic*}, resume, leftmargin=1.6em]
    \item Proposer une analyse physique argumentée de l’influence de ces différents paramètres.
    \item Expliquer pourquoi il n’existe pas nécessairement un \og seul paramètre miracle \fg.
\end{enumerate}

\vspace{0.7em}
\hrule
\vspace{0.5em}

\section*{Barème indicatif}

\begin{itemize}[leftmargin=1.8em]
    \item Exercice 1 : 12 points
    \begin{itemize}
        \item Partie A : 3 points
        \item Partie B : 4 points
        \item Partie C : 2 points
        \item Partie D : 3 points
    \end{itemize}
    \item Exercice 2 : 14 points
    \begin{itemize}
        \item Partie A : 4 points
        \item Partie B : 3 points
        \item Partie C : 4 points
        \item Partie D : 3 points
    \end{itemize}
    \item Exercice 3 : 18 points
    \begin{itemize}
        \item Partie A : 4 points
        \item Partie B : 4 points
        \item Partie C : 6 points
        \item Partie D : 2 points
        \item Partie E : 2 points
    \end{itemize}
\end{itemize}

\end{document}